\section{Introduction}

%\yfc{Why is pagetable verification a necessary task? What is the challenge? Why previous approaches does not work? }

%\yfc{Explain what a page table is. Maybe provide an example of a 3-level. Say what property to verify: behaves like a one-to-one mapping. What is this property? Because the essence of a page table is a mapping from physical address to virtual address, it is natural to model it as an array.}

%Operating systems (OS) use the paging hardware that traverses page tables to isolate memory and control access permissions between running entities. Faulty implementations nullify the intended memory protection. Thus, verifying the functional correctness of page table management code in core systems is essential for ensuring high security assurance. However, this verification is a significant challenge. Page tables are hierarchical, in-memory structures that are combined with complex, hardware-defined bit encoding. The code that manages them relies on sophisticated pointer arithmetic, dereferences, bitwise operations, and is coupled with dynamic page table allocation, making a direct reasoning of its behavior extremely difficult. Previous work~\cite{brun23hotos} stated that page table proofs compose the lion's share of the efforts in formally verifying an operating system.
%\todo{rewrite.}
Modern operating systems rely on paging hardware to virtualize memory and enforce isolation between protection domains. The kernel must construct and maintain hierarchical, in-memory page tables whose layout is constrained by hardware-defined bit encodings, and whose manipulation involves subtle pointer and arithmetic operations. Because even minor bugs in this code can compromise memory safety, it is essential to obtain strong formal guarantees about its correctness. At the same time, the complexity of page-table structures and their low-level manipulation makes such formal verification technically difficult. Prior work~\cite{brun23hotos} reports that verification of low-level system code that manages page tables accounts for the lion's share of the proof efforts.

\begin{comment}
In existing verified systems\cite{li21sp,brun23hotos}, page-table verification is typically divided into two parts: verifying \emph{access-permission properties} and verifying the \emph{correctness of the paging functions} that update and read multi-level page tables. The former is usually straightforward once an abstract view of the page table is available \yfc{give evidence, some papers also said this?}. The latter, i.e., verifying the paging functions themselves, is far more difficult and is the subject of this work. The correctness of paging functions is commonly stated as a \emph{structural lifting} property: the concrete multi-level page table must behave exactly like an abstract flat array mapping virtual addresses to physical ones. 
\end{comment}

%\sw{rewritten by SW.}
Direct verification of the page table management codebase is exceedingly difficult due to the complexity of multi-level tree manipulations. Therefore, prior works~\cite{li:ieeesp2021,brun23hotos,li:sekvm:usenix21} employ a modular proof strategy that first establishes the correctness of the fundamental core paging functions: \texttt{map}, \texttt{unmap}, and \texttt{resolve} (i.e., \texttt{walk}). These functions manage, read, and write to the concrete in-memory multi-level page tables. Past works~\cite{li:ieeesp2021,brun23hotos,li:sekvm:usenix21} express the correctness of the core paging functions as a \textit{structural lifting property}: the concrete multi-level page tables behave exactly like an abstract flat array mapping addresses to page table entry values. 
The correctness of paging functions is established by showing that the concrete implementation satisfies an abstract specification defined over a flat array. This strategy simplifies verification reasoning for higher-level, complex system features, such as copy-on-write or fine-grained access controls, that utilize paging functions. They can then be reasoned over the abstract specifications, substantially reducing proof complexity. 

\emph{Refinement proofs} are one common way to establish the structural lifting property. They show that the concrete code conforms to an abstract specification by proving that every concrete step corresponds to an abstract step that preserves a relation between the two states. This requires reasoning of the form “for all concrete states, there exists a matching abstract state,” which introduces alternating universal–existential quantifiers. Quantifiers of this form are a well-known bottleneck for automated SMT solvers: they can significantly degrade solver performance, cause proof searches to diverge, or lead to solver failure. Consequently, fully automating refinement proofs is difficult in practice.

To cope with these challenges, prior work~\cite{brun23hotos,atmosphere2025sosp,Brun2022VerifiedPaging} that verifies paging functions using SMT-based tools like Verus~\cite{verus} must insert carefully placed \texttt{assert} statements or hand-crafted \emph{triggers} to guide quantifier instantiation. Triggers are syntactic patterns attached to quantified formulas that indicate when the solver may instantiate a quantifier. Although indispensable for controlling SMT reasoning, they are fragile: poorly chosen triggers may cause the solver to miss necessary instantiations, generate irrelevant ones, or fall into non-terminating instantiation loops. These issues become particularly severe for refinement proofs, where the presence of alternating quantifiers forces the solver to coordinate multiple layers of instantiation. Even small changes to the specification or to the memory model can invalidate a previously working trigger strategy. As a result, extensive manual tuning is required, and verification may still fail due to the solver's inherent limitations~\cite{Brun2022VerifiedPaging}. 

This work explores a new approach for proving structural lifting of page tables. As discussed above, the main obstacle is that structural lifting is traditionally expressed through refinement proofs, which introduce quantifiers that SMT solvers handle poorly. Our approach instead captures the core lifting idea in a \emph{quantifier-free} encoding: we verify that concrete paging functions satisfy the \emph{array axiom}, which characterizes arrays in SMT logic. The axiom states that a select-after-store returns the stored value at the updated index and preserves all other indices, i.e.,
\[
\text{Select}(\text{Store}(A,i,v), j) =
\begin{cases}
v & \text{if } i = j,\\
\text{Select}(A,j) & \text{otherwise}.
\end{cases}
\]
We aim to verify that paging functions managing multi-level page tables satisfy this axiom. This verification allows us to abstract the complex operations, interpreting them as updates to a simplified abstract map. This design avoids the quantifier instantiation problems that plague refinement proofs, thereby significantly improving solver robustness and performance.
By eliminating the substantial reasoning burden inherent in refinement-based approaches that require quantifiers, our method makes verification more efficient and scalable, enabling adaptation to diverse page-table implementations.

Based on the proof strategy, we introduce the Paging Function Verification Framework (PPV).
PPV verifies paging-function implementations written in C by applying \emph{symbolic execution} and SMT-based reasoning. To show that PPV can interoperate with existing verification frameworks and to highlight its practical value in enabling sound verification of complex systems, we construct a \emph{connecting proof} integrating PPV with Spoq~\cite{spoq}, an automated framework for full-system verification in Coq. Similar to prior work~\cite{certikos,li:ieeesp2021,taososp21,li:sekvm:usenix21}, PPV focuses on the functional correctness of paging primitives. Concurrency concerns, such as page-table synchronization, TLB shootdowns, and cache maintenance, are handled by separate verified layers and are orthogonal to the structural lifting property considered by PPV. The original Spoq framework requires laborious manual efforts for refinement proofs of paging function correctness. For the SeKVM hypervisor verification~\cite{spoq-artifact}, Spoq relies on 2,500 lines of Coq tactic scripts to prove the \texttt{map}, \texttt{unmap}, and \texttt{resolve} primitives. Integrating PPV enables seamless, push-button replacement of these structural lifting proofs. This flexibility allows system developers to employ alternative designs tailored to specific performance and feature requirements.

%Our infrastructure allows different paging-function implementations to be verified
%automatically with only minimal local adjustments. This flexibility enables
%developers to explore alternative paging designs and choose implementations whose
%page sizes and granularities better match the needs of specific applications. In
%practice, selecting paging primitives with more suitable granularities can improve
%overall system performance. Our results, therefore, show not only that our structural-lifting framework verifies diverse paging mechanisms uniformly but also
%that it can help to enhance the performance of memory management in real systems.

%, \texttt{map}, \texttt{unmap}, and \texttt{resolve}
To demonstrate versatility, we applied PPV to verify multiple C implementations of paging function primitives for the formally verified SeKVM hypervisor that operate on multi-level page tables, support dynamic allocation and multiple mapping granularities, including huge pages. The integration of the PPV proofs with Spoq enables the initial end-to-end correctness verification of SeKVM’s verified codebase to be built on altering paging function implementations.
%To evaluate the integration with commodity low-level systems, we ported the verified implementations to a secure Linux KVM hypervisor~\cite{secvma2024acsac} that manages multi-level page tables.
Our results show that PPV automates the verification of diverse implementations of paging functions with minimal adjustments to the specifications. The push-button approach substantially improves performance efficiency. We evaluated both a refinement-based encoding and our array–axiom–based encoding on a multi-level paging implementation. The refinement proof did not complete even after more than a day, whereas our encoding discharged the proof in roughly ten seconds. %Furthermore, we demonstrated that by leveraging our SMT-based approach, we can upgrade an existing verified paging function implementation~\cite{sekvm-artifact-sosp21} to support an additional 1 GB of mapping granularity while preserving the structural lifting property of page tables, thereby enhancing performance and obviating the need for complex re-verification.
Our work complements a critical gap in the state-of-the-art automated system verification frameworks in proving core low-level systems, such as OSes and hypervisors that manage page tables.

The remainder of this paper is organized as follows. Section~\ref{sec:background} provides background knowledge and discusses related work on the verification of paging functions, together with their limitations. Section~\ref{sec:overview} presents an overview of our quantifier-free proof strategy, introduces the PPV framework, and walks through a detailed case study. Section~\ref{sec:semantics} defines the semantics of our symbolic executor. Section~\ref{sec:eval} evaluates the verification performance of PPV and demonstrates its capability in bug identification and full-system verification. Finally, Section~\ref{sec:concl} offers concluding remarks.

%Our push-button approach substantially improves the efficiency of paging function verification. We evaluated both a refinement-based encoding and our array–axiom–based encoding on a multi-level paging implementation. The refinement proof did not complete even after more than a day, whereas our encoding discharged the proof in roughly ten seconds.  

%Our approach significantly improved the efficiency of the correctness proof of paging functions. We tried to verify a 3-level implementation of paging functions using both refinement proof-based encoding and the array axiom-based ones. The former failed to prove its correctness using more than one day, while the latter proved it in 10 seconds. To demonstrate how versatile our approach is, we applied the it to verify multiple C-based implementations of three fundamental paging primitives—\texttt{map}, \texttt{unmap}, and \texttt{resolve}—that manage multi-level page tables with dynamic allocation and the support of multiple mapping granularities. To evaluate our approach in practical, commodity software, we ported the verified implementations to a secure Linux KVM hypervisor~\cite{secvma2024acsac}.
%Our results demonstrate that the new implementation not only adheres to the correctness properties of structural-lifting but also enhances the performance efficiency of memory management. We believe that our approach, which leverages SMT solvers, can optimize and facilitate the verification of complex, low-level systems, including OS kernels and hypervisors.

%\todo{refine!}
%The practicality of such a specification is also demonstrated by our successful verification of multiple SeKVM-style~\cite{sekvm-artifact-sosp21} page-table implementation variants using this approach. \todo{rewrite this.}
%\todo{further sell our contributions with concrete implementations?}

%\sw{SW: be specific about what we have verified: basic paging functions (or operations), set/get? Or do we call them map/unmap?}

%To address this, we redesigned the formulation of structural lifting proofs to be quantifier-free where possible, making them directly solvable by SMT solvers. We demonstrated its practicality by successfully verifying multiple SeKVM-style~\cite{sekvm-artifact-sosp21} page table implementations.

